\documentclass[12pt]{article}
\usepackage[T1]{fontenc}
\usepackage[utf8]{inputenc}
\usepackage[polish]{babel}
\usepackage{lmodern}
\usepackage[a4paper,margin=2.4cm]{geometry}
\usepackage{xcolor}
\usepackage{graphicx}
\usepackage{amsmath,amssymb}
\usepackage{booktabs}
\usepackage{array}
\usepackage{enumitem}
\usepackage{hyperref}
\usepackage{tcolorbox}
\usepackage{fancyhdr}
\definecolor{WEIIPrimary}{RGB}{0,130,150}
\definecolor{WEIIPrimaryDark}{RGB}{0,86,100}
\definecolor{WEIIPrimaryLight}{RGB}{229,246,248}
\definecolor{PRzNavy}{RGB}{31,63,142}
\definecolor{RLTextDark}{RGB}{34,34,34}
\definecolor{RLGray}{RGB}{90,90,90}
\definecolor{RLLightGray}{RGB}{245,246,248}
\definecolor{RLWarning}{RGB}{184,86,0}
\definecolor{RLExample}{RGB}{52,112,64}
\definecolor{RLCodeBg}{RGB}{248,248,248}

\usepackage{listings}
\usepackage{xcolor}

\lstdefinestyle{rlpython}{%
    language=Python,
    basicstyle=\ttfamily\scriptsize,
    keywordstyle=\color{WEIIPrimaryDark}\bfseries,
    commentstyle=\color{RLGray}\itshape,
    stringstyle=\color{RLExample},
    numbers=left,
    numberstyle=\tiny\color{RLGray},
    numbersep=6pt,
    showstringspaces=false,
    breaklines=true,
    tabsize=4,
    frame=single,
    rulecolor=\color{WEIIPrimary!35},
    backgroundcolor=\color{RLCodeBg},
    captionpos=b
}

\lstset{style=rlpython}


\hypersetup{colorlinks=true,linkcolor=WEIIPrimaryDark,urlcolor=WEIIPrimaryDark,citecolor=WEIIPrimaryDark}
\setlength{\headheight}{15pt}
\pagestyle{fancy}
\fancyhf{}
\lhead{Uczenie się ze wzmocnieniem}
\rhead{WEiI PRz}
\cfoot{\thepage}

\newtcolorbox{taskbox}[1]{colback=WEIIPrimaryLight,colframe=WEIIPrimary,title=\textbf{#1}}
\newtcolorbox{notebox}[1]{colback=RLLightGray,colframe=RLGray,title=\textbf{#1}}
\newtcolorbox{warninghandout}[1]{colback=RLWarning!8,colframe=RLWarning,title=\textbf{#1}}

\newcommand{\handouttitle}[3]{%
\begin{center}
{\Large\bfseries #1}\\[0.2cm]
{\large #2}\\[0.15cm]
{\small #3}
\end{center}
\vspace{0.4cm}
\hrule
\vspace{0.6cm}
}

\newcommand{\EnableLabNumbering}{%
  \renewcommand{\thesection}{L\arabic{section}}%
  \renewcommand{\thesubsection}{\thesection.\arabic{subsection}}%
  \renewcommand{\thesubsubsection}{\thesubsection.\arabic{subsubsection}}%
}


\title{Raport projektowy\\Uczenie się ze wzmocnieniem}
\author{Jakub Jurczyk}
\date{2026/2027}

\begin{document}
\maketitle
\tableofcontents
\newpage

\section{Informacje podstawowe}

\begin{center}
\begin{tabular}{p{0.34\textwidth}p{0.56\textwidth}}
\toprule
Element & Informacja \\
\midrule
Jakub Jurczyk &  \\
179509 &  \\
Informatyka/AI &  \\
Robot-Kurier &  \\
Link do repozytorium & \url{https://github.com/xjjurczyk/UzW-2026-Robot-Kurier/tree/main} \\
Data oddania & Do uzupełnienia \\
\bottomrule
\end{tabular}
\end{center}

\section{Temat i cel projektu}

\subsection{Opis problemu}

W projekcie zajmuję się problemem sterowania robotem-kurierem na dwuwymiarowej planszy. Na mapie umieszczone są przeszkody, punkty ładowania baterii, miejsce odbioru paczki oraz punkt docelowy. Głównym utrudnieniem jest to, że robot posiada ograniczoną baterię.

\subsection{Cel agenta}

Agent ma za zadanie dojść do paczki, podnieść ją i zanieść do punktu docelowego, zanim całkowicie wyczerpie mu się prąd.

\subsection{Uzasadnienie zastosowania uczenia się ze wzmocnieniem}

Należy wyjaśnić, dlaczego problem można rozpatrywać jako problem uczenia się ze wzmocnieniem. W szczególności należy wskazać:
\begin{itemize}
    \item jakie decyzje są podejmowane sekwencyjnie, - Dostarczenie paczki to nie jest jedna prosta decyzja, tylko cała seria kroków na mapie od startu aż do mety.
    \item w jaki sposób decyzje wpływają na przyszły stan środowiska, - Każdy wykonany ruch przesuwa agenta na inne pole i jednocześnie zmniejsza jego zapas energii
    \item dlaczego istotny jest wynik długoterminowy, a nie tylko natychmiastowa nagroda, - Jeśli agent skupi się tylko na najszybszym dojściu do celu, może mu zabraknąć prądu w połowie drogi. Czasem musi zaakceptować tymczasową stratę.
    \item dlaczego zastosowanie agenta uczącego się jest uzasadnione. - RL pozwoli agentowi poprzez próby i błędy samodzielnie nauczyć się balansu między dwoma celami (czas dostawy i oszczędzanie energii).
\end{itemize}

\section{Formalizacja problemu RL/MDP}

\subsection{Agent i środowisko}

Agentem w tym zadaniu jest robot-kurier, a środowiskiem plansza podzielona na kratki. W pojedynczym kroku epizodu agent wybiera jedną z 4 akcji (ruch w górę, dół, lewo lub prawo).

\subsection{Stan lub obserwacja}



\begin{center}
\begin{tabular}{p{0.30\textwidth}p{0.60\textwidth}}
\toprule
Składowa stanu/obserwacji & Znaczenie \\
\midrule
pozycjax & Obecna kolumna, w której znajduje się agent na planszy \\
pozycjay & Obecny wiersz, w którym znajduje się agent na planszy \\
poziomenergii & Stan baterii \\
statuspaczki & Wartość 0 lub 1 (0 – agent szuka paczki, 1 – ma już paczkę i idzie do mety) \\
\bottomrule
\end{tabular}
\end{center}

\subsection{Akcje}



\begin{center}
\begin{tabular}{p{0.30\textwidth}p{0.60\textwidth}}
\toprule
Akcja & Opis \\
\midrule
0 & Ruch w górę \\
1 & Ruch w dół \\
2 & Ruch w lewo \\
3 & Ruch w prawo \\
\bottomrule
\end{tabular}
\end{center}

\subsection{Funkcja nagrody}



\begin{equation}
    r_t = \text{do uzupełnienia}
\end{equation}

W tej sekcji należy wyjaśnić:
\begin{itemize}
    \item za co agent otrzymuje nagrodę dodatnią, - Przyznawana jest za podniesienie paczki (+10) oraz za dostarczenie jej na metę (+100)
    \item za co agent otrzymuje karę, - Agent otrzymuje dużą karę (-100) za rozładowanie baterii lub wpadnięcie w pułapkę.
    \item czy występuje kara za krok, - Występuje stała, niewielka kara za każdy krok (-1)
    \item czy występują ograniczenia twarde lub miękkie, - Ograniczenia twarde nie pozwalają na zmianę pozycji, ale agent i tak dostaje za taki ruch karę (-1)
    \item czy nagroda nie zachęca agenta do niepożądanego zachowania. - Obecny układ (kara za krok -1 i duża kara za śmierć -100) powinien zapobiec sytuacji, w której agent woli od razu zginąć, byleby nie zbierać kolejnych kar za kroki. Wartości te mogą jednak wymagać drobnej korekty podczas eksperymentów.
\end{itemize}

\subsection{Warunki zakończenia epizodu}

Agent z sukcesem dostarczy paczkę do punktu docelowego, Agentowi wyczerpie się energia, Agent wpadnie w pole z pułapką

\subsection{Założenia i ograniczenia środowiska}
Środowisko jest w pełni dyskretne i statyczne (przeszkody nie ruszają się po mapie).
Wybrana akcja zawsze powoduje przemieszczenie np. w górę, o ile na drodze nie ma ściany.
Zużycie energii jest stałe i wynosi dokładnie 1 jednostkę na każdy wykonany ruch.

\section{Implementacja}

\subsection{Struktura repozytorium}

Należy krótko opisać strukturę repozytorium i najważniejsze pliki projektu.

\begin{verbatim}
repozytorium/
|-- src/
|-- experiments/
|-- report/
|-- results/
`-- README.md
\end{verbatim}

\subsection{Środowisko}

Należy opisać sposób implementacji środowiska, w szczególności funkcje inicjalizacji epizodu oraz wykonania kroku środowiska.

\subsection{Agent bazowy}

Należy opisać agenta losowego, heurystycznego albo inny baseline wykorzystany do porównania z agentem uczącym się.

\subsection{Agent uczący się}

Należy opisać zastosowany algorytm RL lub DRL, jego parametry oraz sposób implementacji. W przypadku Q-tab należy opisać reprezentację tablicy wartości. W przypadku DRL należy opisać wejścia modelu, wyjścia modelu, architekturę oraz sposób uczenia.

\section{Plan eksperymentów}

Należy opisać eksperymenty zaplanowane w celu oceny działania agenta.

\begin{center}
\begin{tabular}{p{0.24\textwidth}p{0.30\textwidth}p{0.30\textwidth}}
\toprule
Eksperyment & Co jest zmieniane? & Co jest mierzone? \\
\midrule
E1 & Do uzupełnienia & Do uzupełnienia \\
E2 & Do uzupełnienia & Do uzupełnienia \\
E3 & Do uzupełnienia & Do uzupełnienia \\
\bottomrule
\end{tabular}
\end{center}

W tej sekcji należy wskazać:
\begin{itemize}
    \item liczbę epizodów uczenia,
    \item liczbę powtórzeń eksperymentu albo uzasadnienie braku powtórzeń,
    \item parametry uczenia,
    \item metryki oceny,
    \item warianty środowiska, map lub funkcji nagrody,
    \item sposób porównania z agentem bazowym.
\end{itemize}

\section{Wyniki}

W tej sekcji należy przedstawić wyniki eksperymentów. Wyniki powinny być pokazane w postaci tabel, wykresów lub czytelnych zestawień.

\begin{center}
\begin{tabular}{p{0.28\textwidth}p{0.20\textwidth}p{0.20\textwidth}p{0.20\textwidth}}
\toprule
Wariant & Średni zwrot & Skuteczność & Długość epizodu \\
\midrule
Agent bazowy & Do uzupełnienia & Do uzupełnienia & Do uzupełnienia \\
Agent RL/DRL & Do uzupełnienia & Do uzupełnienia & Do uzupełnienia \\
\bottomrule
\end{tabular}
\end{center}

\begin{figure}[h]
    \centering
    \fbox{\parbox{0.82\textwidth}{\centering Miejsce na wykres procesu uczenia}}
    \caption{Przykładowy wykres procesu uczenia agenta.}
\end{figure}

\section{Analiza i wnioski}

Należy zinterpretować wyniki. W tej części należy odpowiedzieć m.in. na pytania:
\begin{itemize}
    \item Czy agent uczący się działa lepiej niż agent bazowy?
    \item Jak zmieniało się zachowanie agenta w trakcie uczenia?
    \item Które elementy środowiska były najtrudniejsze?
    \item Czy funkcja nagrody prowadziła do oczekiwanego zachowania?
    \item Czy wystąpiły przypadki niepożądanego zachowania agenta?
    \item Jakie są najważniejsze ograniczenia rozwiązania?
\end{itemize}

\section{Instrukcja uruchomienia}

Należy opisać sposób uruchomienia projektu od czystego repozytorium. Przykładowo:

\begin{verbatim}
git clone https://github.com/...
cd nazwa-repozytorium
python -m venv .venv
.venv\Scripts\activate
pip install -r requirements.txt
python src/main.py
\end{verbatim}

Należy również wskazać, jak odtworzyć eksperymenty oraz gdzie zapisywane są wyniki.

\section{Wsparcie AI}

W tej sekcji należy obowiązkowo opisać wykorzystanie narzędzi sztucznej inteligencji. Jeżeli student nie korzystał z takich narzędzi, należy wpisać jednoznaczną informację, że nie korzystano ze wsparcia AI.

\begin{center}
\begin{tabular}{p{0.13\textwidth}p{0.19\textwidth}p{0.21\textwidth}p{0.27\textwidth}}
\toprule
Data & Narzędzie/model & Cel użycia & Transkrypt / link / zakres \\
\midrule
Do uzupełnienia & Do uzupełnienia & Do uzupełnienia & Do uzupełnienia \\
Do uzupełnienia & Do uzupełnienia & Do uzupełnienia & Do uzupełnienia \\
\bottomrule
\end{tabular}
\end{center}

Należy zapisać transkrypty rozmów z modelami AI wraz z datą oraz nazwą modelu. Jeżeli rozmowa odbywała się poza narzędziami programistycznymi, należy dodatkowo udostępnić całą rozmowę poprzez link, o ile użyte narzędzie umożliwia takie udostępnienie.

Student powinien także wyjaśnić, które elementy wygenerowane lub zasugerowane przez AI zostały samodzielnie sprawdzone, zmodyfikowane albo odrzucone.

\section{Literatura}

W tej sekcji należy wskazać wykorzystane źródła: książki, artykuły naukowe, dokumentację bibliotek, materiały wykładowe oraz inne źródła użyte podczas pracy nad projektem.

\begin{thebibliography}{9}
\bibitem{sutton_barto}
R. S. Sutton, A. G. Barto, \textit{Reinforcement Learning: An Introduction}, 2nd ed., MIT Press, 2018.
\end{thebibliography}

\end{document}
